\scanpage{51}
\mspage{\deleted{\unsure{6}}7}

d'une structure déjà orientée si $a$ est une arête libre (correspondant à un $C\in\pi_0(\partial X)$), ou qu'on dispose de l'\deleted{\unsure{une|ens}} des deux orientations $\omega(a)$ \deleted{\unsure{=}} dans le cas où $a$ est une arête ordinaire, implique que l'on tient en mains, pour chaque $a\in A$, un polygone combinatoire $P_a$, dont l'ens des sommets est $R_a$, et l'ens des arêtes peut se décrire \deleted{\unsure{comme}} de façon canonique en termes de $R_a,\omega_a$ et $\omega_a\to\mathrm{Circ}(R_a)$. Il peut être plus commode de la donner d'emblée dans la structure, \deleted{qui serait alors un contour} \inserted{où la famille des $\unsure{R_a|P_a}$ serait \unsure{renforcé|remplacé} par un} \uline{contour combinatoire}, i.e. un ensemble de repères $\underline{\vec R}$ (dits ``repères de rigidification''), avec une opération de
\[
 {\unsure{\mathbb D|D}}_{\infty}=\{\sigma_0,\sigma_1\mid\sigma_0^2=\sigma_1^2=1\}
\]
dessus, et une application
\[
 \underline{\vec R}\longrightarrow\vec A
\]
compatible avec l'hom
\[
 \mathbb D_{\infty}\xrightarrow{\unsure{\varepsilon|2}}\{\pm1\}
\]
sur les groupes d'opérateurs, mais de plus un sous ens
\[
 \underline{\vec R}_{l}^{+}\subset\underline{\vec R}\mid\vec A_l
\]

\margin{($\vec A_l={}$ arêtes libres, correspondant aux cercles de \unsure{recollement|recollements}, i.e. les $\in\pi_0(\deleted{\unsure{\partial}}\partial X)$)}

définissant une orientation du contour correspondant.
% archive sheet number: 50

